\section{Conclusion and Future Work}
\label{sec:conclusion}

This paper presented a situation-aware ISAC co-simulation framework for 6G
smart warehouses. The framework connects Sionna RT with ns-3/5G~LENA so that
site-specific ray tracing, sensing, tracking, adaptive beamforming, and NR
communication can be evaluated in one workflow. The sensing module converts
ray-traced reflections into target detections, the tracking module stabilizes
those detections over time, and the beamforming module uses the confirmed
target positions to update the communication link.

The controlled validation shows that the framework behaves as expected when
the scene becomes more difficult. In free space, visibility is favorable, so
ISAC reduces the mean path loss by up to about $18$~dB and keeps the
false-positive rate low when tracking is enabled. In the warehouse scene, ISAC
still improves propagation, with the widest tested beam giving the strongest
path-loss reduction. However, racks, walls, and partial visibility reduce the
detection rate and increase missed detections. This means that the main
warehouse sensing problem is not pointwise localization accuracy after a target
is detected, but maintaining complete and continuous detections in clutter. The
ablation results also show two important trade-offs: deeper propagation gives
only a small detection gain and can add false positives, while removing the
tracker increases raw detections but makes the output much less clean.

The end-to-end warehouse Digital Twin confirms that these controlled findings
matter at application level. The $2{\times}2$ gNB array is not sufficient for a
stable ISAC gain because robot traffic becomes less reliable. The $4{\times}4$
array gives the best operating point: it improves robot-link propagation and
throughput, keeps the rest of the warehouse traffic usable, improves sensing
stability, and avoids the large power increase of the $8{\times}8$ array. The
$8{\times}8$ case provides more antenna resources, but its extra energy cost is
not matched by a proportional communication or sensing improvement.

These results show that ISAC should not be evaluated from sensing accuracy or
communication throughput alone. A useful warehouse deployment must balance all
three dimensions: reliable target awareness, useful link improvement, and
reasonable energy cost. Future work should therefore focus on more robust
clutter handling, stronger track continuity in dense indoor scenes, and
beam-management methods that combine target position with channel state
information. The evaluation should also be extended with more traffic loads,
random seeds, robot routes, and warehouse layouts so that the operating limits
of the framework can be characterized more completely.
